\newpage
\section*{Термины и определения}
\addcontentsline{toc}{section}{Термины и определения}

\begin{enumerate}

\item \textbf{Виртуализация} --- технология создания программного
(виртуального) представления вычислительных ресурсов (процессора, памяти,
накопителей, сетевых интерфейсов), позволяющая запускать несколько
независимых операционных систем на одном физическом устройстве.

\item \textbf{Виртуальная машина (ВМ)} --- программная эмуляция
вычислительного устройства, предоставляющая гостевой операционной системе
доступ к выделенным ей виртуальным аппаратным ресурсам посредством
виртуализации.

\item \textbf{VirtualBox (Oracle VM VirtualBox)} --- программный продукт
виртуализации для операционных систем Windows, Linux, macOS и других,
позволяющий запускать гостевые операционные системы в изолированной среде
виртуальных машин на базе хост-системы.

\item \textbf{Вычислительный кластер} --- совокупность взаимосвязанных
вычислительных узлов, объединённых высокоскоростной сетью и
функционирующих как единая вычислительная система под управлением
специализированного программного обеспечения.

\item \textbf{IP-адрес (Internet Protocol address)} --- уникальный
числовой идентификатор сетевого интерфейса узла в компьютерной сети,
используемый для адресации и маршрутизации пакетов данных по протоколу IP.

\item \textbf{TCP-соединение (Transmission Control Protocol)} ---
логическое соединение между двумя узлами сети, устанавливаемое по
протоколу TCP и обеспечивающее надёжную, упорядоченную и контролируемую
на предмет ошибок передачу данных между ними.

\item \textbf{SSH (Secure Shell)} --- сетевой протокол прикладного уровня,
позволяющий производить удалённое управление операционной системой и
туннелирование TCP-соединений, в том числе с использованием авторизации по
криптографическим ключам без передачи пароля по сети.

\item \textbf{DNS (Domain Name System)} --- распределённая система
преобразования символьных доменных имён узлов сети в соответствующие им
IP-адреса и обратно.

\item \textbf{Разрешение имён} --- процесс сопоставления символьных имён
узлов сети соответствующим им IP-адресам, выполняемый средствами DNS или
локального файла соответствий (\texttt{/etc/hosts}).

\item \textbf{Диспетчер заданий} --- программный компонент системы
управления кластером, отвечающий за приём заявок от пользователей, ведение
очереди заданий, планирование и распределение ресурсов вычислительных
узлов.

\item \textbf{Управляющий узел (head node, control node)} --- узел
вычислительного кластера, на котором выполняется диспетчер заданий и
средства управления кластером, отвечающие за приём, планирование и
распределение заданий по вычислительным узлам.

\item \textbf{Вычислительный узел (compute node)} --- узел вычислительного
кластера, предоставляющий ресурсы (процессорное время, память) для
непосредственного выполнения заданий, поставленных в очередь диспетчера.

\item \textbf{Munge} --- служба аутентификации, используемая для
проверки подлинности сообщений, которыми обмениваются узлы кластера, на
основе общего секретного ключа.

\item \textbf{SLURM (Simple Linux Utility for Resource Management)} ---
система управления заданиями и ресурсами вычислительного кластера с
открытым исходным кодом, обеспечивающая постановку заданий в очередь, их
планирование и распределение по вычислительным узлам; для аутентификации
узлов использует службу Munge.

\item \textbf{Заявка (задание, job)} --- описание вычислительной задачи,
поставленной пользователем в очередь диспетчера кластера, включающее
запрошенные ресурсы и оценку времени выполнения.

\item \textbf{Очередь заданий диспетчера} --- упорядоченный список заявок,
принятых диспетчером SLURM и ожидающих выделения ресурсов для запуска на
вычислительных узлах.
% TODO: проверить необходимость последних определений
\item \textbf{Среднее арифметическое} --- статистическая характеристика
выборки, равная сумме всех её значений, делённой на их количество.

\item \textbf{Дисперсия} --- статистическая характеристика разброса
значений случайной величины, равная среднему арифметическому квадратов
отклонений значений от среднего арифметического выборки.

\item \textbf{Стандартное отклонение} --- статистическая характеристика
разброса значений случайной величины относительно её среднего
арифметического, равная квадратному корню из дисперсии выборки.

\item \textbf{Нормальное распределение} --- непрерывное распределение
вероятностей случайной величины, полностью определяемое математическим
ожиданием и стандартным отклонением; график плотности имеет симметричную
колоколообразную форму.

\item \textbf{Логнормальное распределение} --- распределение вероятностей
случайной величины, логарифм которой имеет нормальное распределение с
параметрами $\mu$ и $\sigma$.

\end{enumerate}